\documentclass[11pt,a4paper]{article}
\usepackage[flat]{pinns-course}

\coursemeta{Session @N@, computational lab --- TITLE}{Session @N@, computational lab}{M2 MACIA, AGM, CY Cergy Paris University}

% ================================================================
\begin{document}

\coursetitle{Session @N@ \textbullet\ Computational lab}{TITLE}{}

\vspace{6pt}

\begin{plainbox}
\textbf{What this is.} A short companion to SESSION. N tasks.

\smallskip
\textbf{Material.} \code{lab@N@-student.ipynb} and \code{lab@N@-solutions.ipynb}.
\end{plainbox}

\vspace{4pt}

\begin{warnbox}
\textbf{The rule of this lab.} \emph{Predict before you run.} A computation that
confirms what you already proved teaches you nothing; one that contradicts what you
expected teaches you a great deal.
\end{warnbox}

% ================================================================
\section{Task 1 --- PLACEHOLDER}

\begin{plainbox}
\textbf{Recalls.}
\end{plainbox}

\begin{task}[TITLE \tagcore]\label{task:a}
\leavevmode
\begin{enumerate}[label=\textbf{\arabic*.}]
\item
\end{enumerate}
\end{task}

% One task must establish something no theorem of the session addresses.
% If no task does, the lab is optional. See docs/authoring.md, §5.

% ================================================================
\ifsolutions
\newpage
\section{Expected results, and what to say about them}
\addcontentsline{toc}{section}{Expected results}

% Every number here must come from an actual run of lab/experiments.py.
\fi

\end{document}
